\documentclass[11pt,a4paper]{article}
\usepackage[T1]{fontenc}
\usepackage[utf8]{inputenc}
\usepackage[polish]{babel}
\usepackage{lmodern}
\usepackage{geometry}
\usepackage{booktabs}
\usepackage{siunitx}
\usepackage{pgfplots}
\usepackage{pgfplotstable}
\usepackage{csvsimple}
\usepackage{underscore}
\usepackage{hyperref}
\usepackage{caption}
\usepackage{float}

\pgfplotsset{compat=1.18}
\geometry{margin=2.2cm}
\sisetup{group-separator={\,}}

\newcommand{\DataDir}{l2-benchmarks/data}

\title{\textbf{UniVerify: Benchmarki warstwy 2\\w rejestrze dyplomów opartym na Ethereum}\\[0.6em]
\large Porównanie Sepolia, Arbitrum, Base oraz zkSync}
\author{Nazar Shcherbyna \\ Praca magisterska}
\date{Pomiar wykonano: 2026-05-12T21:11:58.253Z\unskip}

\begin{document}
\maketitle
\tableofcontents
\newpage

\section{Wprowadzenie}
Niniejszy rozdział porównuje koszt uruchomienia kontraktu \texttt{DiplomaRegistry}
na czterech sieciach Ethereum: warstwie pierwszej (Sepolia, jako punkt odniesienia)
oraz trzech rozwiązaniach warstwy drugiej reprezentujących różne rodziny
(Arbitrum Nitro, Base jako przedstawiciel OP-stack, zkSync Era jako rollup z dowodami zwięzłości).

\section{Metodologia}
\subsection{Pomiary}
Każda sieć mierzona jest na sieci testowej: \emph{zużycie gazu} i
\emph{opłata za dane L1} są wielkościami rzeczywistymi,
natomiast koszty w USD są modelowane z~historycznych wartości \texttt{baseFeePerGas}
warstwy pierwszej z ostatnich 90 dni (percentyle p10/p50/p90).

\subsection{Parametry próby}
Dla \texttt{issueBatch} użyto rozmiarów paczek
$\{1, 10, 100, 1\,000, 10\,000\}$. Dla \texttt{revokeFromBatch}
zebrano 30 próbek przeciwko paczce seed o rozmiarze 40.
Dla \texttt{statusWithProof} zebrano 100 próbek czasu odpowiedzi RPC.

\section{Wyniki: koszt \texttt{issueBatch}}
\begin{table}[H]
  \centering
  \caption{Tabela 1: Koszt \texttt{issueBatch} według sieci i rozmiaru paczki.}
  \csvautotabular{\DataDir/table-issue-cost.csv}
\end{table}

\begin{table}[H]
  \centering
  \caption{Tabela 2: Koszt na dyplom (amortyzacja paczki) w USD, scenariusz p50.}
  \csvautotabular{\DataDir/table-issue-per-diploma.csv}
\end{table}

\begin{figure}[H]
  \centering
  \begin{tikzpicture}
    \begin{loglogaxis}[
      xlabel={Rozmiar paczki (liczba dyplomów)},
      ylabel={Koszt na dyplom (USD)},
      legend pos=north east,
      grid=major,
      width=0.85\textwidth,
      height=8cm,
    ]
      \addplot table[x=batchSize, y=sepolia] {\DataDir/fig1-cost-per-diploma.dat};
      \addplot table[x=batchSize, y=arbitrumSepolia] {\DataDir/fig1-cost-per-diploma.dat};
      \addplot table[x=batchSize, y=baseSepolia] {\DataDir/fig1-cost-per-diploma.dat};
      \addplot table[x=batchSize, y=zksyncSepolia] {\DataDir/fig1-cost-per-diploma.dat};
      \legend{Sepolia L1, Arbitrum Sepolia, Base Sepolia, zkSync Sepolia}
    \end{loglogaxis}
  \end{tikzpicture}
  \caption{Rys. 1: Koszt na dyplom w funkcji rozmiaru paczki (skala log-log, scenariusz p50).}
\end{figure}

\section{Wyniki: koszt \texttt{revokeFromBatch}}
\begin{table}[H]
  \centering
  \caption{Tabela 3: Koszt \texttt{revokeFromBatch} (pojedyncza transakcja, mediana).}
  \csvautotabular{\DataDir/table-revoke-cost.csv}
\end{table}

\section{Wyniki: latencja}
\begin{table}[H]
  \centering
  \caption{Tabela 4: Latencja odczytu \texttt{statusWithProof}.}
  \csvautotabular{\DataDir/table-read-latency.csv}
\end{table}

\begin{table}[H]
  \centering
  \caption{Tabela 5: Czas odbioru paragonu transakcji zapisu po stronie klienta.}
  \csvautotabular{\DataDir/table-inclusion-latency.csv}
\end{table}

\noindent Mierzona wartość to czas między wywołaniem \texttt{sendTransaction}
a otrzymaniem paragonu przez klienta RPC. Przy interwale odpytywania
około 4 s (domyślnym w bibliotece viem) wynik jest górnym ograniczeniem
faktycznego czasu inkluzji w bloku, nie jego dokładnym pomiarem;
rozkład wartości na sieci Sepolia jest wyraźnie skupiony wokół
wielokrotności czasu bloku ($\approx 12$ s), co potwierdza powyższą
interpretację.

\section{Wyniki: wrażliwość na scenariusze basefee}
\begin{figure}[H]
  \centering
  \begin{tikzpicture}
    \begin{axis}[
      ybar,
      symbolic x coords={sepolia, arbitrumSepolia, baseSepolia, zksyncSepolia},
      xtick=data,
      ylabel={Koszt \texttt{issueBatch(1000)} (USD)},
      bar width=10pt,
      legend pos=north west,
      width=0.85\textwidth,
      height=8cm,
    ]
      \addplot table[x=chain, y=quiet_usd] {\DataDir/fig4-basefee-scenarios.dat};
      \addplot table[x=chain, y=normal_usd] {\DataDir/fig4-basefee-scenarios.dat};
      \addplot table[x=chain, y=congested_usd] {\DataDir/fig4-basefee-scenarios.dat};
      \legend{p10 (quiet), p50 (normal), p90 (congested)}
    \end{axis}
  \end{tikzpicture}
  \caption{Rys. 4: Koszt \texttt{issueBatch(1000)} przy scenariuszach basefee p10/p50/p90.}
\end{figure}

\section{Dyskusja}

Wszystkie pomiary opisane w niniejszym rozdziale wykonano przez jednego dostawcę RPC w jednej lokalizacji POP, z $N{=}30$ powtórzeniami dla każdego pomiaru opóźnienia oraz losowym wyborem liści przy revoke. Wyniki to średnia $\pm\sigma$ z trzech niezależnych runów. Wszystkie parametry konfiguracyjne (ziarno losowości, dostawca, region, kurs ETH/USD) są zapisane w pliku \texttt{docs/l2-benchmarks/data/meta.json}.

\subsection{Ograniczenia pomiaru}

Poniżej zebrano metodologiczne ograniczenia, których czytelnik powinien być świadomy interpretując wyniki. Cztery z pierwotnie zidentyfikowanych ograniczeń zostały wyeliminowane w toku rewizji eksperymentu (jednolity dostawca RPC, $N{=}30$ powtórzeń, losowy wybór liści do revoke, trzy niezależne runy z raportowaniem $\sigma$); pozostałe sześć omówiono poniżej.

\paragraph{Założenie kompresji Brotli dla Arbitrum.} Model kosztu zapisu danych na L1 dla Arbitrum przyjmuje współczynnik kompresji Brotli równy 1.0, ponieważ dane Merkle (hashe pseudolosowe) nie kompresują się znacząco. Realny aplikacyjny calldata o strukturze rzadkiej kompresuje się 2--4-krotnie, co oznacza, że nasze oszacowanie kosztu L1 dla Arbitrum jest \textbf{górnym ograniczeniem}; rzeczywiste koszty produkcyjne mogą być proporcjonalnie niższe.

\paragraph{Niewspółmierność jednostki gazu w zkSync Era.} Wartość \texttt{gasUsed} raportowana przez zkSync (131\,463 dla \texttt{issueBatch} z $n{=}1$) łączy w jednej jednostce koszt wykonania, dane publiczne i narzut bootloadera kontekstu Account Abstraction --- nie jest to ta sama jednostka, co \texttt{gasUsed} w EVM. Porównanie surowych liczb gazu między rodzinami rollupów jest \textbf{metodologicznie niepoprawne}; sensowne porównanie kosztów może odbywać się wyłącznie w USD lub ETH.

\paragraph{Opóźnienie inkluzji jest ograniczone przez interwał odpytywania.} Pomiar czasu od submit do receipt opiera się o cykl odpytywania klienta viem (\textasciitilde 4\,s). Otrzymane wartości są zatem \textbf{górnym ograniczeniem czasu inkluzji w bloku}, a nie samym czasem inkluzji. Porównanie międzysieciowe jest dodatkowo zniekształcone interakcją czasu bloku z interwałem odpytywania.

\paragraph{Snapshot skalarów Base SystemConfig.} Model kosztu Base Ecotone używa snapshotu parametrów \texttt{baseFeeScalar=2269} oraz \texttt{blobBaseFeeScalar=1055762} pobranego z bloku L1 \texttt{25087416}. Po zmianie tych skalarów przez operatora Base wynikowe koszty zmieniłyby się proporcjonalnie; data snapshotu jest udokumentowana w \texttt{meta.json} każdego eksportu.

\paragraph{Snapshot kursu ETH/USD.} Wyniki w USD przeliczono kursem $\text{ETH}/\text{USD} = 3500$ (snapshot CoinGecko z dnia 2026-05-13). Wynik dla innego kursu można uzyskać mnożeniem przez czynnik $\frac{\text{nowy kurs}}{3500}$.

\paragraph{Konserwatywne ceny sekwensera L2.} Ceny gazu sekwensera L2 zostały ustawione na zaokrąglone wartości powyżej żywego snapshotu (Arbitrum 0.1\,gwei vs 0.02 żywe; Base 0.005 vs 0.006; zkSync 0.05 vs 0.0453), aby pochłonąć typowe wahania kongestii.

\subsection{Reprodukowalność}

Wszystkie wyniki w tym rozdziale można odtworzyć z czystego klona repozytorium UniVerify poleceniem \texttt{npm run benchmarks:all -{}- -{}-runs=3} po wypełnieniu pliku \texttt{.env} kluczem prywatnym portfela testowego oraz adresami RPC czterech sieci testnet (Sepolia, Arbitrum Sepolia, Base Sepolia, zkSync Sepolia). Pipeline wykonuje sekwencję \texttt{deploy} $\to$ \texttt{measure} $\to$ \texttt{price} $\to$ \texttt{aggregate} $\to$ \texttt{export}, regenerując wszystkie tabele i wykresy w \texttt{docs/l2-benchmarks/data/}. Każdy artefakt jest powiązany z plikiem \texttt{meta.json} zawierającym \texttt{runId}, \texttt{measuredAt}, kurs ETH/USD z datą snapshotu, ziarno losowości (\texttt{BENCH\_RANDOM\_SEED}), dostawcę i region RPC oraz listę łańcuchów. Plik \texttt{packages/benchmarks/README.md} dokumentuje wymagania środowiskowe i znane ograniczenia. CI repozytorium (\texttt{.github/workflows/benchmarks-drift.yml}) automatycznie regeneruje tabele z zacommitowanych próbek wyników i sygnalizuje rozjazd między źródłem a zatwierdzonymi danymi.

\paragraph{Wniosek.}
\dots{}

\end{document}
